% Chinese build uses ctexart (xelatex): it auto-configures a system CJK font and
% renders reliably, whereas the pdfLaTeX-oriented AIAA class breaks xeCJK routing.
\documentclass[a4paper,10pt,twocolumn]{ctexart}
\usepackage[margin=2cm]{geometry}
\usepackage{graphicx}
\usepackage{amsmath}
\usepackage{siunitx}
\usepackage{longtable,tabularx}
\usepackage{booktabs}
\usepackage{float}
\usepackage[section]{placeins}
\usepackage{xcolor}
\usepackage{array}
\usepackage{xurl}
\usepackage[numbers,sort&compress]{natbib}
\bibliographystyle{new-aiaa}
\graphicspath{{figures/}}
\emergencystretch=2.5em
\AtBeginDocument{%
  \setlength{\intextsep}{5pt plus 2pt minus 1pt}%
  \setlength{\textfloatsep}{6pt plus 2pt minus 2pt}%
  \setlength{\abovecaptionskip}{3pt}%
  \setlength{\belowcaptionskip}{1pt}%
}

\title{单斜置机翼跨介质飞行器的 A 层几何规划尺寸设计：\\CFD 标定代理模型与受浮力限制的载荷前沿}
\author{张一杨\thanks{独立研究者，温彻斯特公学，英国温彻斯特。}}
\date{}

\begin{document}
\maketitle

\section{方法}

\subsection{建模框架与 A 层的范围}
本飞行器为一细长的导弹形艇体，其单一直机翼安装在一个中央竖直枢轴上。在空中，机翼垂直展
开于艇身以产生升力；在入水之前，机翼旋转至与艇身平行贴合，因而在水下飞行器成为一个洁净的
回转体，仅以其少量残余外伸段提供水动力配平力。对这样一个飞行器进行尺寸设计，需要耦合三个
通常彼此独立的学科——空气动力学、水动力学与结构——以及两种动压相差两个数量级以上的工作介
质。为使这一耦合问题保持可解，并且——尤为关键地——为保证所求得的最优解是\emph{全局}最优，
本设计被分解为三层。A 层，即本文的主题，是一个几何规划（GP）：它在给定的机翼斜置角与给定
的浮力状态下，确定全部连续尺寸变量。B 层是一个在下潜过程中调度机翼斜置轨迹的贝叶斯优化器，
位于 A 层之上，本文不作讨论。一项计算流体力学（CFD）研究为这两层提供无法用解析式写出的空
气动力与水动力系数。

几何规划是这样一类优化问题：在正项式（posynomial）不等式约束与单项式（monomial）等式约束
下，最小化一个正项式目标。\emph{单项式}是设计变量若干次幂之积，且前置系数为正；\emph{正项
式}则是若干单项式之和。使该表述具有吸引力的决定性性质在于：在变量替换 $u_i\mapsto\ln u_i$
之下，每个 GP 都成为凸问题，因而求解器返回的是全局最优——而非仅仅局部最优~\cite{boyd2007tutorial}——
同时给出目标对每一约束的灵敏度。这一保证的工程代价，是每一条控制方程都必须被写成正项式或
单项式形式；下文的方法，在很大程度上正是诚实地支付这一代价的过程。

\subsection{CFD 试验设计}
A 层无法从第一性原理导出的系数为：空中与水下的阻力极曲线、升力线斜率及其随机翼斜置的变化、
收拢机翼的配平升力，以及空化吸力峰。正是这些量——且仅有这些量——界定了 CFD 试验设计。两项
研究在 OpenFOAM~v2606 中以定常不可压求解器 \texttt{simpleFoam} 及 $k$--$\omega$ SST 湍流模
型完成。空中研究在一个 $17\times12\times12$~m 的自由空气域内，对从收拢（$\Lambda=90^\circ$）
到展开（$\Lambda=0^\circ$）、每 $10^\circ$ 一档的全部十个机翼位置进行了网格划分，并通过旋转
入口速度矢量而非重新划分网格，将迎角从 $-4^\circ$ 扫掠至 $+14^\circ$，由此节省了约九成的网
格划分工作量。水下研究复用同一网格并改用水的物性，在展开与收拢两极加入了 $2$ 至 $20$~m/s
的速度扫掠，并将工作深度作为后处理参数处理——因为在不可压解中流场与绝对压强水平无关，深度
仅通过空化数进入。所有力系数均一致地以机翼平面面积为参考，收敛以系数稳定性判定：要求最后
五十次迭代的均值与前五十次之差小于百分之一，而非以固定的残差阈值判定。作为对整条流程的独
立校核，收拢机翼的诱导阻力因子在空气与水中——雷诺数相差约 $1.5$ 倍——一致到 $7\%$ 之内；由
于阻力系数无量纲、且在雷诺数匹配时仅依赖于几何，这一一致性证明两项研究彼此自洽。

\subsection{GP 兼容的代理模型拟合}
每条 CFD 关系均以 \texttt{gpfit} 流程~\cite{hoburg2016data} 化为 GP 合法的代理式。凡在对数
坐标下呈仿射的关系，均以对数空间的普通最小二乘拟合为单一单项式——这在数学上等同于单项最大
仿射拟合，且数值稳健；凡带有真实曲率的关系，则改以多项软极大仿射（softmax-affine）正项式拟
合。升力线斜率随机翼斜置的变化为
\begin{equation}
C_{L\alpha}(\Lambda) = 0.847\,(\cos\Lambda)^{1.105}\ \text{rad}^{-1},
\qquad(\text{rms}_{\log}=0.011),
\label{eq:clalpha}
\end{equation}
其指数 $1.105$ 与斜置无关性原理所预言的 $1$ 相差不超过百分之十一，故该拟合不止于插值：它印
证了“只有垂直于机翼的速度分量产生升力”这一解析预期。诱导阻力因子随机翼折叠、其有效展弦比
坍缩而急剧上升，
\begin{equation}
k(\Lambda) = 0.631\,(\cos\Lambda)^{-1.352}\ ,
\qquad(\text{rms}_{\log}=0.013),
\label{eq:kschedule}
\end{equation}
自完全展开时的 $0.64$ 攀升至收拢时的 $22.9$——这三十六倍的增幅，正是关系式 $k=1/(\pi e A)$
中展弦比坍缩的直接标志。水下零升阻力随雷诺数标度为 $C_{D0,w}=0.185\,\mathrm{Re}^{-0.112}$，
其指数落在纯湍流摩擦阻力的 $-0.2$ 与纯形状阻力的 $0$ 之间，与一个阻力由二者混合构成的物体
相符。空化吸力峰被发现基本为常数——作为无量纲量本应如此：对于收拢构型，其配平升力使之接近
零升力系数，故 $-C_{p,\min}\approx1.16$，此常数即用作保守的空化门限。

\subsection{A 层几何规划}
目标是最大化载荷，表述为最小化 $m_{\text{pay}}^{-1}$。约束集由随附的完整表述文档转录而来，
其书写方式使每一条关系都是被一个单项式所界定的正项式，因而在 GP 中是可接纳的。空中巡航、
阻力分解与浮力—配平各块为
\begin{align}
mg &\le \tfrac12\rho_a V_a^2 S\,C_{L,a}, & C_{L,a} &\le C_{L,\max}, \label{eq:aeriallift}\\
C_{D,a} &\ge \frac{C_{L,a}^2}{\pi e A} + C_{f}(1+k)\frac{S_{\text{wet}}}{S} + C_{Dp,a},
      & C_f &\ge \frac{0.074}{\mathrm{Re}^{0.2}}, \label{eq:aerialdrag}\\
(1+k) &\ge 1 + 1.5\!\left(\tfrac{d}{L}\right)^{3/2} + 7\!\left(\tfrac{d}{L}\right)^{3},
      & B &= \rho_w g\,\nabla, \label{eq:formfactor}\\
mg &\le B\,(1+\xi), & L_{\text{tr}} &\le \tfrac12\rho_w V_w^2 S\,C_{L,\text{uw}}. \label{eq:buoytrim}
\end{align}
式~\eqref{eq:aeriallift} 中第一条是 Hoburg 与 Abbeel~\cite{hoburg2014geometric} 的定常飞行升
力平衡：巡航时展开机翼须承担全部重量，且因升力系数越大机翼面积越小，故写为一条被优化器驱至
等号的不等式。它与一条取自 CFD 失速数据的硬上限 $C_{L,a}\le C_{L,\max}$ 相配，否则优化器将
索求物理上不可能的升力系数以无止境地缩小机翼。式~\eqref{eq:aerialdrag} 是经典的阻力分
解~\cite{hoburg2014geometric,hoerner1965fluid}，分为诱导项、零升黏性项与机翼翼型型阻项。诱
导项 $C_{L,a}^2/(\pi e A)$ 为升力线理论结果~\cite{hoerner1985fluid}；奥斯瓦尔德效率取解析先
验值 $e\approx0.85$ 而非受网格限制的 CFD 值，理由见第~\ref{sec:justify} 节。零升项
$C_f(1+k)S_{\text{wet}}/S$ 以平板摩擦系数乘以形状因子与湿面积—参考面积之比；此处采用湍流幂
律 $C_f=0.074\,\mathrm{Re}^{-0.2}$ 以代替 ITTC~1957 对数律，因其为单项式且在相关雷诺数范围
内误差小于百分之二~\cite{hoburg2014geometric,myring1976theoretical}。式~\eqref{eq:formfactor}
为 Hoerner 对流线型回转体的形状因子~\cite{hoerner1965fluid}，其中 $1.5(d/L)^{3/2}$ 项为超速
增量、$7(d/L)^3$ 项为分离—压差分量；对我们的细长比 $d/L\approx0.10$，其值为
$(1+k)\approx1.054$，即一个百分之五的黏性附加量，远低于早期手算所假设的百分之二十，故我们
以此取代那一保守估计。式~\eqref{eq:formfactor} 中的浮力为对排开艇体体积 $\nabla$ 的阿基米德
定律，而式~\eqref{eq:buoytrim} 中的竖直平衡 $mg\le B(1+\xi)$ 表明重量由浮力加上由配平承担的
分量 $\xi$ 所支撑；该配平力由收拢机翼的少量外伸段提供，
$L_{\text{tr}}\le\tfrac12\rho_w V_w^2 S\,C_{L,\text{uw}}$，这正是自补偿论点的解析支柱——折叠机
翼虽展弦比坍缩，仍能在不付出大阻力代价的情况下配平飞行器。水下阻力、空化、能量、结构与可行
性门限各块为
\begin{align}
C_{D,w} &\ge C_{D0,w} + k_w\,C_{L,\text{uw}}^2, & D_w &\ge \tfrac12\rho_w V_w^2 S\,C_{D,w},\label{eq:hydrodrag}\\
\tfrac12\rho_w V_w^2\,(-C_{p,\min}) &\le p_{\text{atm}} + \rho_w g h - p_v,
      & E_b &\ge P_a\,t_a + P_w\,t_w, \label{eq:cavenergy}\\
t_s &\ge \frac{\rho_w g h\,d}{2\sigma_y}, & \mathrm{SF}\cdot\rho_w g h &\le 2.42\,E\!\left(\tfrac{t_s}{d}\right)^{2.5}\!\!\frac{d}{L}, \label{eq:structure}\\
S_{\text{fin}}\,l_{\text{fin}} &\ge V_{H,\min}\,S\,\bar c, & C_{l,\text{dist}} &\le C_{L\alpha,\text{fin}}\,\delta_{\max}\frac{S_{\text{fin}}}{S}\frac{b}{L}. \label{eq:gates}
\end{align}
式~\eqref{eq:hydrodrag} 是收拢构型的水下阻力极曲线，其形式与空中极曲线相同（诱导加寄生），
但诱导因子 $k_w$ 直接取自水动力研究，因为收拢机翼的展弦比坍缩是任何解析先验都无法干净刻画
的真实效应；与模型中每个系数一样，阻力以机翼平面面积为参考，使 CFD 数据无须改变参考即可转
入。式~\eqref{eq:cavenergy} 含两条约束。其一为 Brennen~\cite{brennen2014cavitation} 的空化起
始条件 $-C_{p,\min}\le\sigma$，经整理后要求水面处的吸力压强不得低于饱和蒸气压；Amromin 与
Rozhdestvensky~\cite{amromin2022cavitation} 指出，在我们约 $2\times10^{7}$ 量级的艇体雷诺数
下，最小压力系数与起始数之差趋于消失，故等式 $\sigma_i=-C_{p,\min}$ 不仅保守、且愈发精确，
该约束由此对水下速度设下一个稳固的上限。其二为两段航程的能量平衡：对电池电驱动飞行器，燃烧
式飞机的燃油消耗模型被恒质量平衡取代，其中储存能量须覆盖空中段与水下段的推进功~\cite{allen2000propulsion,lin2026energies}，
电池质量则由该能量除以电芯级比能得出。式~\eqref{eq:structure} 以两个极限中较严者确定耐压壳
壁厚：单项式的薄壳环向应力量规，以及外压屈曲坍缩——对后者，将信号式（signomial）的 Bryant
表达式以 Windenburg--Trilling 闭式单项式加二倍安全系数取代，遵循针对小直径无加强筋管件的推
荐做法~\cite{shinoka2024structural}。式~\eqref{eq:gates} 施加两道可行性门限。静稳定性——其精
确形式含 Renilson~\cite{renilson2018submarine} 竖直面内符号相异的导数、因而为信号式——通过
尾容量系数 $V_H=S_{\text{fin}}l_{\text{fin}}/(S\bar c)$~\cite{hoerner1985fluid} 施加；要求其
超过一常规下限，是保证足够静稳定裕度的单项式代理。滚转操纵权要求十字形尾翼在偏转至极限时，
产生的滚转力矩系数至少等于空中研究所测得的斜置诱导机翼力矩峰值，采用
Ref.~\cite{hoerner1985fluid} 的低展弦比尾翼升力斜率。入水生存被有意排除于此约束集之外、并于
其后单独施加，理由见第~\ref{sec:justify} 节。最后，扫掠 $\xi$ 使式~\eqref{eq:buoytrim} 在每
次求解时保持单项式，故整个程序始终为纯 GP；若放开 $\xi$，该等式将成为信号式而须改用序贯求解
器~\cite{hoburg2014geometric,boyd2007tutorial}。

\section{模型论证}\label{sec:justify}

\subsection{为何采用几何规划，以及为何扫掠浮力状态}
选择几何规划不是出于便利，而是出于可信。因对数变换使问题成为凸问题，求解器可确证全局最优，
并额外、无偿地报告每条约束对目标的推力大小；在同一非凸问题上使用基于梯度的优化器，则二者皆
不可得。其代价——将每条关系写成正项式——在此处是负担得起的，因为控制物理本就接近幂律形式：
阻力按雷诺数的幂标度、浮力为几何尺寸之积、诱导阻力为升力的平方。

浮力失衡 $\xi=(W-B)/B$ 度量飞行器重量偏离其所排开水重的程度。$\xi=0$ 时飞行器中性浮、悬停而
无须配平力；$\xi>0$ 时偏重，收拢机翼须产生向上的配平升力以保持深度；$\xi<0$ 时偏轻，机翼向
下配平。将失衡作为参数扫掠而非放开，出于一个明确的理由：$\xi$ 固定时，平衡 $mg=B(1+\xi)$ 为
单项式，整个模型保持为纯 GP；放开 $\xi$ 将使该等式成为信号式，须改用较慢的序贯（信号式规
划）求解器。扫掠 $\xi$ 也是物理上更具信息量的选择，因为它直接描出载荷—浮力前沿：更重的飞行
器有更大的质量预算、因而载荷更多，直到机翼在空中再也无法举起它为止。

\subsection{采用解析阻力并以 CFD 标定，而非固定的 CFD 极曲线}
将实测的展开机翼阻力极曲线直接代入 GP 是诱人的。此法被否决，有两个原因。其一，固定极曲线绑
定于单一几何；一旦艇体与机翼成为设计变量（第~\ref{sec:findings} 节），阻力就必须随之标度，
而这只有解析分解——摩擦系数乘形状因子乘湿面积比——才能做到。其二，也更微妙的是，由扫掠网格
推得的展开机翼效率并未收敛：拟合所得升力线斜率 $0.835$~rad$^{-1}$ 约为该机翼展弦比下薄翼预
期的五分之一，而拟合所得诱导因子意味着奥斯瓦尔德效率仅 $0.07$，远低于预期的 $0.8$--$0.9$。
两处偏差指向同一根源——一个 $57{,}000$ 单元的扫掠网格对机翼随迎角的响应分辨不足——故其绝对
效率取自解析先验 $e\approx0.85$，胜于取自受网格限制的拟合。CFD \emph{确实}可靠给出的，是
式~\eqref{eq:clalpha}--\eqref{eq:kschedule} 的无量纲\emph{标度关系}、收拢极曲线（其展弦比坍
缩是真实效应而非网格伪象）、空化门限，以及经交叉验证的形状因子；模型所用的正是这些。

\subsection{固定与自由几何、可行性门限，以及排除入水}
模型在两种构型下求解。第一种将几何固定于既有 CAD 尺寸，用以核验现有飞行器能否满足任务；第
二种在一个发射包线内放开艇体直径与长度、机翼平面形状及尾翼面积，用以揭示优化器所偏好的设
计。静稳定性要求——其精确形式含符号相异的稳定性导数、因而为信号式——通过标准尾容量系数
$V_H$ 与 $V_V$ 施加；要求二者超过常规下限，是一个 GP 合法的代理，可在不离开凸世界的前提下保
证相同的静裕度行为。最后，入水生存被排除于优化之外，有两个原因：峰值减速系数尚待其自身的二
维冲击研究，且冲击约束是一个通过/失败的生存判据，并不像任务能量与浮力约束那样对飞行器起定
尺作用。因此它在求解之后作为对已定型设计的审计来评估，而非作为其中的驱动因素。

\section{结果}\label{sec:findings}

\subsection{CFD 拟合系数}
拟合系数在物理上自洽，且凡有解析先验之处均能予以还原。图~\ref{fig:polars} 给出全部十个斜置
角的阻力极曲线：诱导阻力因子随机翼折叠平滑而单调地增大，无任何异常点，且每条按机翼分列的极
曲线都以高于 $0.98$ 的决定系数被复现。图~\ref{fig:clalpha} 绘出升力线斜率随斜置的变化，并叠
加式~\eqref{eq:clalpha} 的单项式拟合；接近于一的指数印证了 $\cos\Lambda$ 律。图~\ref{fig:kcollapse}
在对数坐标上给出诱导因子的展弦比坍缩，而图~\ref{fig:cavitation} 说明空化数据为何分为两支——
收拢构型处于接近零升处，而展开构型跨越一段真实升力区间——因而如最初尝试那样将二者拟合为一
条曲线是不正确的；收拢支给出门限值。

\begin{figure}[h]\centering
\includegraphics[width=0.86\linewidth]{aero_polars_all.png}
\caption{全部十个机翼斜置角的空中阻力极曲线。诱导阻力因子 $k$（图例）自完全展开
（$\Lambda=0^\circ$）的 $0.64$ 单调升至收拢（$\Lambda=90^\circ$）的 $22.9$。}
\label{fig:polars}
\end{figure}

\begin{figure}[h]\centering
\includegraphics[width=0.72\linewidth]{sched_CLalpha.png}
\caption{升力线斜率随机翼斜置的变化。拟合指数 $1.105$ 将 $C_{L\alpha}\propto\cos\Lambda$ 的斜
置无关性先验印证到百分之十一之内。}
\label{fig:clalpha}
\end{figure}

\begin{figure}[h]\centering
\includegraphics[width=0.72\linewidth]{sched_k.png}
\caption{对数坐标下诱导阻力因子随机翼斜置的变化：三十六倍的增幅是机翼折平时展弦比坍缩的标
志。}
\label{fig:kcollapse}
\end{figure}

\begin{figure}[h]\centering
\includegraphics[width=0.80\linewidth]{cavitation_split.png}
\caption{空化吸力峰随升力系数的变化。收拢构型（接近零升）与展开构型（有限升力区间）形成两条
不同的支；收拢支 $-C_{p,\min}\approx1.16$ 设定空化门限。}
\label{fig:cavitation}
\end{figure}

\subsection{载荷—浮力前沿}
图~\ref{fig:payload} 对两种几何情形描出载荷随浮力失衡的变化。在两种情形中，载荷均随 $\xi$ 上
升，因更重的飞行器有更大的质量预算；在两种情形中，前沿都终止于空中升力系数达到其失速上限、
机翼在空中再也无法支撑飞行器之处。固定几何飞行器在 $\xi$ 的可行范围内承载 $2.4$ 至 $3.5$~kg
载荷，这证实既有设计满足数千克量级的载荷目标。放开几何使载荷大致翻倍，至 $5.9$ 至 $8.4$~kg，
其原因由第~\ref{sec:optgeom} 节精确给出：优化器将艇体扩大至发射包线边缘，从而“买到”浮力。

\begin{figure}[h]\centering
\includegraphics[width=0.72\linewidth]{payload_xi.png}
\caption{固定 CAD 几何与受包线约束的自由几何下，载荷随浮力失衡 $\xi$ 的变化。两条曲线均终止
于空中升力系数达到 $0.60$ 失速上限之处。}
\label{fig:payload}
\end{figure}

\subsection{优化后的几何及其约束}\label{sec:optgeom}
表~\ref{tab:soln} 列出跨失衡扫掠的自由几何解。其模式一致而富有启示：在每个 $\xi$ 值上，艇体
直径与长度都顶到 $0.12$~m 与 $1.30$~m 的包线上限，机翼展长顶到“一个艇长减 $0.30$~m”的几何
上限，空中升力系数顶到其 $0.60$ 失速上限。换言之，优化器同时将飞行器推向三个极限——发射平台
允许它多大、机身允许机翼展多宽、机翼在空中能出多大力——而它返回的载荷，正是在支付了任务能量
与耐压壳之后所剩者。由静稳定性门限定尺的十字形尾翼，总面积约 $86$~cm$^2$、力臂 $0.585$~m，增
重约四分之三千克；滚转操纵权门限则被稳定性本已要求的尾翼富余满足。两道门限都不驱动载荷，这
恰是可行性门限应有的表现。

\begin{table}[h]\centering
\caption{跨浮力扫掠的自由几何 A 层解。$\xi=0.15$ 一行为参考设计点。}\label{tab:soln}
\small
\begin{tabular}{@{}rrrrrrrr@{}}\toprule
$\xi$ & 载荷 & 质量 & $d$ & $L$ & 展长 & $S$ & $S_{\text{fin}}$ \\
      & (kg) & (kg) & (mm) & (mm) & (m) & (m$^2$) & (cm$^2$) \\\midrule
$0.00$ & $5.86$ & $9.12$ & $120$ & $1300$ & $1.00$ & $0.097$ & $64.8$ \\
$0.05$ & $6.28$ & $9.57$ & $120$ & $1300$ & $1.00$ & $0.102$ & $71.4$ \\
$0.10$ & $6.71$ & $10.03$ & $120$ & $1300$ & $1.00$ & $0.107$ & $78.4$ \\
$0.15$ & $7.14$ & $10.48$ & $120$ & $1300$ & $1.00$ & $0.112$ & $85.7$ \\
$0.20$ & $7.56$ & $10.94$ & $120$ & $1300$ & $1.00$ & $0.117$ & $93.3$ \\
$0.25$ & $7.99$ & $11.39$ & $120$ & $1300$ & $1.00$ & $0.122$ & $101.2$ \\
$0.30$ & $8.41$ & $11.85$ & $120$ & $1300$ & $1.00$ & $0.127$ & $109.5$ \\\bottomrule
\end{tabular}
\end{table}

\subsection{灵敏度与设计杠杆}
对偶解为各杠杆排序。最强者为排开体积，灵敏度 $-1.43$：扩大艇体是增加载荷最直接的途径，这也
是优化器将艇体顶到边界的原因。其次为电池比能 $-0.49$，再次为空中效率与航程约 $\pm0.30$，以
及所需水下速度 $+0.38$。收拢诱导阻力因子的灵敏度仅 $+0.001$——可忽略——因为收拢机翼在几乎零
升下工作、其配平阻力也相应极小。这最后一个数字值得驻足，因为它是飞行器设计论点的量化形式：
折叠机翼几乎不付出阻力代价，故相较于入水前抛弃机翼，保留机翼所带来的载荷惩罚很小。

\subsection{代理模型与已定型模型的验证}
每个代理式与已装配的程序，在被信任之前，都经过五种独立手段的校核。其一为前已述及的空—水交叉
验证：收拢诱导阻力因子在两项研究之间一致到百分之七，而这是任何单项研究都无法做出的比较。其
二为与解析先验的吻合：式~\eqref{eq:clalpha} 中的升力斜率指数 $1.105$ 对比于 $1$、被摩擦与形
阻两极所夹的雷诺指数 $-0.112$、以及空化峰接近于零的指数，都落在第一性原理所指之处。其三为复
现质量：每条按机翼分列的极曲线都以高于 $0.98$ 的决定系数被还原，每个保留的单项式拟合都以低
于 $0.02$ 的对数均方根误差被还原。其四为 GP 合法性，逐项验证：进入约束的每个量都具有正系数
形式，唯一的例外——有弯度的展开机翼、其零升截距为负——被一个正的软极大仿射代理式取代。其五
为物理单调性：升力随迎角上升、阻力随机翼展开上升、吸力峰与速度无关，在所采样的每一点上皆
然。已装配的程序以同样的精神受审计：其对偶灵敏度的符号与量级均如预期，其起作用的约束——质量
闭合与浮力——正是一个接近中性的跨介质飞行器所应受限者，这本身就是模型接线正确的证据。

\section{A 层结论}

A 层研究得出三条结论。第一，飞行器受浮力与包线限制，而非受空气动力或结构限制：起作用的约束
是艇体尺寸上限与空中失速上限，而主导灵敏度指向排开体积，故载荷几乎是直接以发射平台所允许的
尺寸“买来”的。既有 CAD 飞行器承载约三千克载荷；同族中经包线优化的飞行器承载约七千克。第
二，设计论点经受住了数字的检验：收拢机翼的配平阻力在尺寸设计中可忽略，故相较抛弃机翼，将一
副可操纵机翼保留至入水几乎不施加载荷惩罚，而升力线斜率洁净的 $\cos\Lambda$ 标度也印证了斜置
机翼的行为符合经典后掠理论所预言者。第三，模型对其最薄弱之处是诚实的：展开机翼的绝对效率尚
未网格收敛，因而经由解析先验而非受网格限制的拟合来承载；实测升力区间上的阻力极曲线过窄，不
足以支撑一个完整的软极大仿射代理；入水常数则有待其自身的研究。这三项正是一项精细化研究应瞄
准的输入——展开机翼的网格收敛网格、更宽的迎角扫掠或翼型级型阻极曲线、以及二维冲击研究——此
后同一 A 层程序在结构不变的情况下将返回更精确的数字。当前模型的价值在于：它已经告诉设计者应
往何处看——依次为发射包线的大小、电池的能量密度，以及展开机翼的效率。

% \bibliographystyle 由 new-aiaa.cls 设定
\bibliography{uauv}

\end{document}
